\documentclass[../linear-algebra.tex]{subfiles}

\begin{document}
\section{Inverse of a Matrix}
\begin{figure}[H]
    \centering
    \includegraphics[width=0.9\textwidth]{matrix-mul-ab.png} % Replace 'filename' with your image name
    \caption*{Transforming by B first then A is written as AB}
    \label{fig:span geometric interpretation}
\end{figure}

If $A$ is an invertible $n\times n$ matrix, then for each $\mathbf{x}$
in $\mathbb{R}^n$, the equation $A\mathbf{x} = \mathbf{b}$
has the unique solution \[\mathbf{x} = A^{-1}\mathbf{b}\]
\textit{Proof}
\begin{enumerate}
    \item Show that solution exists \newline
    Substitute $A^{-1}\mathbf{b}$ into $A\mathbf{x}=\mathbf{b}$, \newline
    LHS = $A\mathbf{x}=A(A^{-1})\mathbf{b}=(AA^{-1})\mathbf{b}=I\mathbf{b}=\mathbf{b}$ = RHS.
    \item Show that solution is unique. \newline
    Suppose there is some vector $u$ such that $A\mathbf{u}=\mathbf{b}$, \newline
    then $A^{-1}A\mathbf{u}=A^{-1}\mathbf{b} \Rightarrow I\mathbf{u}=A^{-1}\mathbf{b} \Rightarrow \mathbf{u}=A^{-1}\mathbf{b}$.
\end{enumerate}

\bigskip

\noindent \textbf{Property}: If $A$ is invertible, then $A^{-1}$ is invertible, and 
\[(A^{-1})^{-1}=A\]

\bigskip

\noindent \textbf{Property}: If $A$ and $B$ are $n\times n$ invertible matrices,
then so is $AB$ and 
\[(AB)^{-1}=B^{-1}A^{-1}\]

\noindent\textit{Proof} \newline
To find some matrix $C$ s.t. $(AB)C=I$ and $C(AB)=I.$ \newline
If $C=B^{-1}A^{-1}$, then $AB(B^{-1}A^{-1})=A(BB^{-1})A^{-1}=AIA^{-1}=AA^{-1}=I$. \newline
Similar argument applies for $(B^{-1}A^{-1})(AB)=I$.

\bigskip

\noindent \textbf{Property}: If $A$ is an invertible matrix, then so is $A^T$, and 
\[(A^T)^{-1}=(A^{-1})^T\]
\textit{Proof} \newline
To show that $(A^{-1})^T$ is the inverse of $A^T$. \newline
Considering the product of $A^T$ and $(A^{-1})^T$,
\[A^T(A^{-1})^T=(A^{-1}A)^T=I^T=I\]
and the product of $(A^{-1})^T$ and $A^T$,
\[(A^{-1})^TA^T=(AA^{-1})^T=I^T=I\]

\subsection{Elementary matrix}
An elementary matrix $E$ is created by performing exactly one elementary row operation
on an identity matrix $I$. \newline
Multiplying a matrix $A$ by $E$ on the left ($EA$)
performs that same row operation on $A$.

\noindent For example,
\[E = \begin{bmatrix}
    0 & 1 & 0 \\
    1 & 0 & 0 \\
    0 & 0 & 1
\end{bmatrix} \text{ and }A=\begin{bmatrix}
    a & b & c \\
    d & e & f \\
    g & h & i
\end{bmatrix}\Rightarrow EA=\begin{bmatrix}
    d & e & f \\
    a & b & c \\
    g & h & i
\end{bmatrix}\]
Every elementary matrix is invertible, and the inverse is another
elementary matrix that reverses the original transformation.
For example,
\[E_1=\begin{bmatrix}
    0 & 1 & 0 \\
    1 & 0 & 0 \\
    0 & 0 & 1
\end{bmatrix} \Rightarrow E_1^{-1}=\begin{bmatrix}
    0 & 1 & 0 \\
    1 & 0 & 0 \\
    0 & 0 & 1
\end{bmatrix}\]
\[E_2=\begin{bmatrix}
    1 & 0 & 0 \\
    0 & 1 & 0 \\
    -4 & 0 & 1
\end{bmatrix} \Rightarrow E_2^{-1}=\begin{bmatrix}
    1 & 0 & 0 \\
    0 & 1 & 0 \\
    4 & 0 & 1
\end{bmatrix}\]

\bigskip

\noindent An $n\times n$ matrix $A$ is invertible iff $A$ is row equivalent
to $I_n$. The same sequence of EROs that reduces $A$ to $I_n$ transforms $I_n$ to $A^{-1}$.

\noindent\textit{Proof} \newline
Suppose $A$ is an invertible $n\times n$ matrix, then $A\mathbf{x}=\mathbf{b}$ has a unique solution for any $\mathbf{b}$.
This means the RREF form of $A$ has a pivot in every row and column. For a $n\times n$ matrix, this has to be $I_n$.
Therefore $A$ is row equivalent to $I_n$. \newline
There exists some sequence of EROs that transforms $A$ to $I_n$, and each row operation
corresponds to multiplication by an elementary matrix. This can be expressed as
\[(E_k\cdots E_2E_1)A=I_n\]
$A$ is invertible since RHS=$I_n$ is invertible and the product of elementary matrices is invertible.\newline
Multiplying by $A^{-1}$ on the right on both sides,
\[(E_k\cdots E_2E_1)=A^{-1}\]
Therefore,
\[(E_k\cdots E_2E_1)I_n=A^{-1}\]

\subsubsection{Example}
To find the inverse of $A = \begin{bmatrix} 0 & 1 & 2 \\ 1 & 0 & 3 \\ 4 & -3 & 8 \end{bmatrix}$, we augment it with the identity matrix $I$ and perform Gauss-Jordan elimination:

\[
\left[ \begin{array}{ccc|ccc}
0 & 1 & 2 & 1 & 0 & 0 \\
1 & 0 & 3 & 0 & 1 & 0 \\
4 & -3 & 8 & 0 & 0 & 1
\end{array} \right]
\]

\textbf{Step 1: Interchange $R_1$ and $R_2$ ($R_1 \leftrightarrow R_2$)}
\[
\xrightarrow{R_1 \leftrightarrow R_2}
\left[ \begin{array}{ccc|ccc}
1 & 0 & 3 & 0 & 1 & 0 \\
0 & 1 & 2 & 1 & 0 & 0 \\
4 & -3 & 8 & 0 & 0 & 1
\end{array} \right]
\]

\textbf{Step 2: Eliminate $a_{31}$ ($R_3 \leftarrow R_3 - 4R_1$)}
\[
\xrightarrow{R_3 - 4R_1}
\left[ \begin{array}{ccc|ccc}
1 & 0 & 3 & 0 & 1 & 0 \\
0 & 1 & 2 & 1 & 0 & 0 \\
0 & -3 & -4 & 0 & -4 & 1
\end{array} \right]
\]

\textbf{Step 3: Eliminate $a_{32}$ ($R_3 \leftarrow R_3 + 3R_2$)}
\[
\xrightarrow{R_3 + 3R_2}
\left[ \begin{array}{ccc|ccc}
1 & 0 & 3 & 0 & 1 & 0 \\
0 & 1 & 2 & 1 & 0 & 0 \\
0 & 0 & 2 & 3 & -4 & 1
\end{array} \right]
\]

\textbf{Step 4: Scale $R_3$ to get a leading 1 ($R_3 \leftarrow \frac{1}{2}R_3$)}
\[
\xrightarrow{\frac{1}{2}R_3}
\left[ \begin{array}{ccc|ccc}
1 & 0 & 3 & 0 & 1 & 0 \\
0 & 1 & 2 & 1 & 0 & 0 \\
0 & 0 & 1 & 1.5 & -2 & 0.5
\end{array} \right]
\]

\textbf{Step 5: Eliminate $a_{13}$ and $a_{23}$ ($R_1 \leftarrow R_1 - 3R_3$ and $R_2 \leftarrow R_2 - 2R_3$)}
\[
\xrightarrow{\substack{R_1 - 3R_3 \\ R_2 - 2R_3}}
\left[ \begin{array}{ccc|ccc}
1 & 0 & 0 & -4.5 & 7 & -1.5 \\
0 & 1 & 0 & -2 & 4 & -1 \\
0 & 0 & 1 & 1.5 & -2 & 0.5
\end{array} \right]
\]

Therefore, the inverse matrix is
\[
A^{-1} = \begin{bmatrix}
-4.5 & 7 & -1.5 \\
-2 & 4 & -1 \\
1.5 & -2 & 0.5
\end{bmatrix}
\]

\newpage
% \begin{minipage}{\textwidth}
\subsubsection{Invertible Matrix Theorem}
Let $A$ be an $n\times n$ matrix. Then the following statements are logically equivalent:
\begin{enumerate}
    \item $A$ is an invertible (non-singular) matrix.
    \item $A$ is row equivalent to $I_n$.
    \item $A$ has $n$ pivot positions.
    \item $A\mathbf{x}=\mathbf{0}$ has only the trivial solution.
    \item The columns of $A$ form a linearly independent set.
    \item $A\mathbf{x}=\mathbf{b}$ has at least one (in fact, unique) solution for each $\mathbf{b}$ in $\mathbb{R}^n$.
    \item The columns of $A$ span $\mathbb{R}^n$.
    \item There is an $n\times n$ matrix $C$ such that $CA=I$.
    \item There is an $n\times n$ matrix $D$ such that $AD=I$.
    \item $A^T$ is an invertible matrix.
    \item The columns of $A$ form a basis for $\mathbb{R}^n$.
    \item $\mathbf{C}(A) = \mathbb{R}^n$.
    \item $\dim \mathbf{C}(A) = n$.
    \item Rank of $A = n$.
    \item $\mathbf{N}(A) = \{\mathbf{0}\}$.
    \item $\dim \mathbf{N}(A) = 0$.
\end{enumerate}


\begin{figure}[H]
    \centering
    \includegraphics[width=0.6\textwidth]{images/invertible-matrix-thm-2.png} % Replace 'filename' with your image name
    \caption*{The links between the statements of invertible matrix theorem.}
    \label{fig:span geometric interpretation}
\end{figure}

\noindent 4: $A\mathbf{x}=\mathbf{0}\Rightarrow\mathbf{x}=A^{-1}\mathbf{0}$.

\noindent $8\Rightarrow 4$: Given $C$ such that $CA=I$,
\[CA\mathbf{x}=C\mathbf{0}\Rightarrow I\mathbf{x}=\mathbf{0}\Rightarrow\mathbf{x}=\mathbf{0}\]
$9\Rightarrow 6$: Given $D$ such that $AD=I$, choose $\mathbf{x}=D\mathbf{b}$,
\[A\mathbf{x}=AD\mathbf{b}=I\mathbf{b}=\mathbf{b}\]
% \end{minipage}

\newpage
\section{Matrix Factorisation}
To solve sequence of equations like $Ax=b_1, Ax=b_2,\ldots, Ax=b_p$,
it is inefficient to first compute the inverse $A^{-1}$ and then multiply it by each vector, $A^{-1}b_{1}, ..., A^{-1}b_{p}$.

\noindent Instead, the efficient solution is to decompose $A_{m \times n}$ into $L_{m \times m}U_{m \times n}$.

\[
L = \begin{bmatrix} 
1 & 0 & \cdots & 0 \\ 
l_{21} & 1 & \cdots & 0 \\ 
\vdots & \vdots & \ddots & \vdots \\ 
l_{m1} & l_{m2} & \cdots & 1 
\end{bmatrix},
\qquad
U = \begin{bmatrix} 
u_{11} & u_{12} & \cdots & u_{1n} \\ 
0 & u_{22} & \cdots & u_{2n} \\ 
\vdots & \vdots & \ddots & \vdots \\ 
0 & 0 & \cdots & u_{mn} 
\end{bmatrix}
\]
where $L$ is unit lower triangular matrix, $U$ is upper triangular matrix.

\subsection{Deriving L \& U}
\textbf{Assumption} $A$ can be reduced to echelon form without requiring any row interchanges.

\noindent First reduce matrix $A$ to $U$, using elementary row operations.
Multiply $A$ by a sequence of unit lower triangular elementary matrices ($E_{1}, \dots, E_{p}$), such that 

\begin{align*}
           & E_{p} \dots E_{1}A = U \\
  \implies & A = (E_{p} \dots E_{1})^{-1}U
\end{align*}
Then $L = E_{p} \dots E_{1}$.

\noindent\textbf{Example} Consider a $3 \times 3$ matrix:
\[
A = \begin{bmatrix} \mathbf{2} & -2 & 3 \\ \mathbf{6} & -7 & 14 \\ \mathbf{4} & -8 & 30 \end{bmatrix}
\]

\noindent Step 1: Row Reduction to find $U$.
Apply elementary row operations to eliminate the entries below the pivots to reach echelon form. 

\noindent First Column: The first pivot is the $\mathbf{2}$ in row 1, column 1. Eliminate the entries below it.
Multiply by $E_{21}$ (which represents $R_2 - 3R_1$):
\[
E_{21}A = \begin{bmatrix} 2 & -2 & 3 \\ 0 & \mathbf{-1} & 5 \\ 4 & -8 & 30 \end{bmatrix}
\]
Multiply by $E_{31}$ (which represents $R_3 - 2R_1$):
\[
E_{31}(E_{21}A) = \begin{bmatrix} 2 & -2 & 3 \\ 0 & \mathbf{-1} & 5 \\ 0 & \mathbf{-4} & 24 \end{bmatrix}
\]

\noindent Second Column: The new pivot is the $\mathbf{-1}$ in row 2, column 2. Eliminate the entry below it.
    * Multiply by $E_{32}$ (which represents $R_3 - 4R_2$):
\[
E_{32}(E_{31}E_{21}A) = \begin{bmatrix} 2 & -2 & 3 \\ 0 & -1 & 5 \\ 0 & 0 & 4 \end{bmatrix} = U
\]

\noindent Step 2: Two Methods to Construct $L$

\noindent Since $A$ has 3 rows, $L$ will be a $3 \times 3$ unit lower triangular matrix. 

\noindent\textbf{Method 1}: The "Multiplier" Method (Fastest).

\noindent At each pivot column, divide the eliminated entries by the pivot for that column, and place the result directly into $L$. 

\noindent Column 1 Multipliers: The pivot is $\mathbf{2}$. Divide 6 and 4 by 2 respectively.

\noindent Column 2 Multipliers: The new pivot is $\mathbf{-1}$. Divide -4 by -1.

\noindent Placing these multipliers below the diagonal of 1s gives you $L$:
\[
L = \begin{bmatrix} 1 & 0 & 0 \\ \mathbf{3} & 1 & 0 \\ \mathbf{2} & \mathbf{4} & 1 \end{bmatrix}
\]

\noindent\textbf{Method 2}: The Elementary Inverse Method.
Mathematically, $L$ is formed by the inverse of the elementary matrices used to reach $U$: $L = E_{21}^{-1}E_{31}^{-1}E_{32}^{-1}$.

\noindent Construct $L$ by simply taking the nonzero off-diagonal elements of the elementary inverse matrices and placing them into the appropriate positions in $L$. 

\[
L = \begin{bmatrix} 1 & 0 & 0 \\ \mathbf{3} & 1 & 0 \\ 0 & 0 & 1 \end{bmatrix} \begin{bmatrix} 1 & 0 & 0 \\ 0 & 1 & 0 \\ \mathbf{2} & 0 & 1 \end{bmatrix} \begin{bmatrix} 1 & 0 & 0 \\ 0 & 1 & 0 \\ 0 & \mathbf{4} & 1 \end{bmatrix} = \begin{bmatrix} 1 & 0 & 0 \\ \mathbf{3} & 1 & 0 \\ \mathbf{2} & \mathbf{4} & 1 \end{bmatrix}
\]

\subsection{Solving using $L$ \& $U$}
\begin{align*}
           & Ax = b \\
  \implies & LUx = b
\end{align*}
\noindent\textbf{Forward substitution}
Defining $y= Ux$, solve for
\[Ly=b\]

\noindent\textbf{Backward substitution}
Solve
\[Ux=y\]

\noindent\textbf{Example}

\noindent Given the linear system $Ax = b$, where:
\[
A = \begin{bmatrix} 3 & -7 & -2 & 2 \\ -3 & 5 & 1 & 0 \\ 6 & -4 & 0 & -5 \\ -9 & 5 & -5 & 12 \end{bmatrix}, \qquad
b = \begin{bmatrix} -9 \\ 5 \\ 7 \\ 11 \end{bmatrix}
\]

\noindent Matrix $A$ has been factored into $L$ and $U$:
\[
L = \begin{bmatrix} 1 & 0 & 0 & 0 \\ -1 & 1 & 0 & 0 \\ 2 & -5 & 1 & 0 \\ -3 & 8 & 3 & 1 \end{bmatrix}, \qquad
U = \begin{bmatrix} 3 & -7 & -2 & 2 \\ 0 & -2 & -1 & 2 \\ 0 & 0 & -1 & 1 \\ 0 & 0 & 0 & -1 \end{bmatrix}
\]

\vspace{0.5cm}
\noindent\textbf{Forward Substitution ($Ly = b$)}

\noindent Set up the lower triangular system to solve for the intermediate vector $y$:
\[
\begin{bmatrix} 1 & 0 & 0 & 0 \\ -1 & 1 & 0 & 0 \\ 2 & -5 & 1 & 0 \\ -3 & 8 & 3 & 1 \end{bmatrix}
\begin{bmatrix} y_1 \\ y_2 \\ y_3 \\ y_4 \end{bmatrix}
= \begin{bmatrix} -9 \\ 5 \\ 7 \\ 11 \end{bmatrix}
\]

\noindent Solve the equations sequentially from top to bottom:
\[
\begin{aligned}
y_1 &= -9 \\
-y_1 + y_2 &= 5 &&\implies y_2 = -4 \\
2y_1 - 5y_2 + y_3 &= 7 &&\implies y_3 = 5 \\
-3y_1 + 8y_2 + 3y_3 + y_4 &= 11 &&\implies y_4 = 1
\end{aligned}
\]

\noindent This gives us our intermediate vector $y$:
\[
y = \begin{bmatrix} -9 \\ -4 \\ 5 \\ 1 \end{bmatrix}
\]

\vspace{0.5cm}
\noindent\textbf{Step 2: Backward Substitution ($Ux = y$)}

\noindent Set up the upper triangular system using the newly found vector $y$ to solve for $x$:
\[
\begin{bmatrix} 3 & -7 & -2 & 2 \\ 0 & -2 & -1 & 2 \\ 0 & 0 & -1 & 1 \\ 0 & 0 & 0 & -1 \end{bmatrix}
\begin{bmatrix} x_1 \\ x_2 \\ x_3 \\ x_4 \end{bmatrix}
= \begin{bmatrix} -9 \\ -4 \\ 5 \\ 1 \end{bmatrix}
\]

\noindent Solve the equations sequentially from bottom to top:
\[
\begin{aligned}
-x_4 &= 1 &&\implies x_4 = -1 \\
-x_3 + x_4 &= 5 &&\implies x_3 = -6 \\
-2x_2 - x_3 + 2x_4 &= -4 &&\implies x_2 = 4 \\
3x_1 - 7x_2 - 2x_3 + 2x_4 &= -9 &&\implies x_1 = 3
\end{aligned}
\]

% \vspace{0.5cm}
\noindent\textbf{Final Solution:}
\[
x = \begin{bmatrix} 3 \\ 4 \\ -6 \\ -1 \end{bmatrix}
\]

\subsection{Zero Pivots: $PA = LU$ Factorization}
The standard $A = LU$ factorization assumes matrix $A$ can be reduced to echelon form without row interchanges. If a pivot becomes zero, you must swap rows, which breaks the standard algorithm. 

\noindent To solve this, we pre-swap the rows of $A$ using a Permutation Matrix $P$, which is an Identity matrix with its rows swapped. This yields the generalized factorization:
\[
PA = LU
\]

\noindent\textbf{Example:}
Given matrix $A$:
\[
A = \begin{bmatrix} 1 & 1 & 1 \\ 1 & 1 & -1 \\ 2 & -1 & -1 \end{bmatrix}
\]

\noindent If we attempt standard row reduction ($R_2 - R_1$ and $R_3 - 2R_1$), we hit a zero pivot in the second column:
\[
\begin{bmatrix} 1 & 1 & 1 \\ 0 & \mathbf{0} & -2 \\ 0 & -3 & -3 \end{bmatrix}
\]

\noindent To fix this, we swap Row 2 and Row 3 using permutation matrix $P_{23}$:
\[
P = \begin{bmatrix} 1 & 0 & 0 \\ 0 & 0 & 1 \\ 0 & 1 & 0 \end{bmatrix}
\]

\noindent Calculate $PA$ to generate the row-swapped matrix:
\[
PA = \begin{bmatrix} 1 & 0 & 0 \\ 0 & 0 & 1 \\ 0 & 1 & 0 \end{bmatrix} \begin{bmatrix} 1 & 1 & 1 \\ 1 & 1 & -1 \\ 2 & -1 & -1 \end{bmatrix} 
= \begin{bmatrix} 1 & 1 & 1 \\ 2 & -1 & -1 \\ 1 & 1 & -1 \end{bmatrix}
\]

\noindent Now, factor $PA$ normally into $L$ and $U$:
Using pivot 1 in column 1: $R_2 - \mathbf{2}R_1 \implies l_{21} = 2$, and $R_3 - \mathbf{1}R_1 \implies l_{31} = 1$.
The entry below the second pivot ($-3$) is already $0 \implies l_{32} = 0$.

\noindent This gives our final $L$ and $U$ matrices for $PA$:
\[
L = \begin{bmatrix} 1 & 0 & 0 \\ 2 & 1 & 0 \\ 1 & 0 & 1 \end{bmatrix}, \qquad
U = \begin{bmatrix} 1 & 1 & 1 \\ 0 & -3 & -3 \\ 0 & 0 & -2 \end{bmatrix}
\]


\end{document}